\emph{The ring.} The function $f$ is bilinear over $\mathbb{F}_{2}$ in the indicator vectors of $S$ and $T$ (each term counts pairs, and a pair counted in $S \triangle S'$ is counted in exactly one of $S$, $S'$ modulo $2$). Hence $f(S \triangle T, U) + f(S, T) \equiv f(S, U) + f(T, U) + f(S, T) \equiv f(S, T \triangle U) + f(T, U) \pmod 2$, which says $(e_{S}e_{T})e_{U} = e_{S}(e_{T}e_{U})$; extended bilinearly the product is associative, with unit $e_{\emptyset}$, and the structure constants are $\pm 1$, so $\Cl_{\Z}(p,q)$ is a ring and, the same computation holding over any commutative ring $R$, $\Cl_{R}(p,q) = R \otimes_{\Z}\Cl_{\Z}(p,q)$ is an $R$-algebra with $R$ central. The rule is the one Proposition \claimref{std-algebra.orthogonal-anticommute} gives over the reals, with $\prod_{i \in S \cap T}Q(e_{i}) = (-1)^{|S \cap T \cap N|}$; it is used here as a comparison, the ring being defined by the rule. For $v = \sum a_{i}e_{i}$, $v^{2} = \sum_{i}a_{i}^{2}e_{i}^{2} + \sum_{i < j}a_{i}a_{j}(e_{i}e_{j} + e_{j}e_{i}) = Q(v)$, since $e_{j}e_{i} = -e_{i}e_{j}$ for $i \ne j$ ($f(\{j\}, \{i\}) - f(\{i\}, \{j\}) = \pm 1$) and $e_{i}^{2} = (-1)^{[i \in N]}$. The check verifies the product rule against GA.py for every pair of blades in all signatures with $n \le 4$, the bilinearity on all triples, and $v^{2} = Q(v)$ on a box of integer vectors.

\emph{(i)} If $Q(v)$ is a unit of $R$, put $u = Q(v)^{-1}v \in \Cl_{R}(p,q)$; then $vu = uv = Q(v)^{-1}v^{2} = 1$, the formula of Proposition \claimref{std-algebra.inverse-of-vector}, here proved over $R$. Conversely, if $vw = 1$ for some $w$, then $v^{2}w^{2} = v(vw)w = vw = 1$, that is $Q(v)\,w^{2} = 1$ since $v^{2} = Q(v)1$ is central; if $wv = 1$, likewise $w^{2}v^{2} = 1$. Comparing grade-$0$ coefficients in the free $R$-module, $Q(v)\,\langle w^{2}\rangle_{0} = 1$, so $Q(v)$ is a unit of $R$, and the one-sided inverse equals $u$ by associativity ($w = (uv)w = u(vw) = u$, or $w = w(vu) = (wv)u = u$). The check verifies, for all integer vectors with coordinates in $\{-3, \dots, 3\}$ and $Q(v) \ne 0$ in $\Cl(2,0)$, $\Cl(1,1)$, $\Cl(0,3)$, $\Cl(1,2)$, that $v/Q(v)$ is a two-sided inverse and that its coefficients lie in $\Z[1/P]$ exactly when $Q(v)$ is a unit there, for $P = \emptyset, \{2\}, \{3\}, \{2,3\}, \{2,5\}$, counts the vectors invertible over $\Z$ against a hand count ($Q(v) = \pm 1$), and, over the ring $\Z/6\Z$ (not a subring of $\Q$), finds by exhaustive search in $\Cl(1,1)$ that a vector has an inverse exactly when $Q(v)$ is prime to $6$.

\emph{(ii)} If every nonzero integer is a unit of $R$, every $v \in V_{\Z}$ with $Q(v) \ne 0$ is invertible by (i). Conversely, if every such $v$ is invertible in $\Cl_{R}(p,q)$, take $v = me_{1}$, $m \ge 1$: $Q(v) = m^{2}Q(e_{1}) = \pm m^{2} \ne 0$, so by (i) $m^{2}$ is a unit of $R$, and so is $m$ ($m \cdot m\,(m^{2})^{-1} = 1$). A commutative ring in which every nonzero integer is a unit receives exactly one ring homomorphism from $\Q$, $a/m \mapsto a\,m^{-1}$, so it is a $\Q$-algebra and $\Q$ is initial among such rings; $\Cl_{\Q}(p,q) = \Q \otimes_{\Z} \Cl_{\Z}(p,q)$ is the localisation of $\Cl_{\Z}(p,q)$ at the nonzero integers, which are central. A subring of $\Q$ that is a $\Q$-algebra contains $1/m$ for all $m$, so it is $\Q$. Explicitly $(me_{1})^{-1} = me_{1}/(m^{2}Q(e_{1})) = e_{1}/(mQ(e_{1}))$ and $e_{1}(me_{1})^{-1} = e_{1}^{2}/(mQ(e_{1})) = 1/m$, so $a/m = (ae_{1})(me_{1})^{-1}$. A vector of $V_{\Q}$ with $Q(v) \ne 0$ is invertible over $\Q$ by (i). The grade-$0$ part of the free $\Q$-module $\Cl_{\Q}(p,q)$ is $\Q e_{\emptyset}$. The check computes $e_{1}(me_{1})^{-1} = 1/m$ and $(ae_{1})(me_{1})^{-1} = a/m$ for $m = 1, \dots, 30$ in three signatures, and the primes that must be inverted, $2, 3, 5, \dots, 29$.

\emph{(iii)} If $p, q \ge 1$, $(e_{1} + e_{p+1})^{2} = 1 - 1 = 0$. If $p = 0$ or $q = 0$, $Q$ is plus or minus a sum of squares on $V_{\Q}$ and vanishes only at $v = 0$. A nonzero $v$ with $v^{2} = 0$ is a zero divisor. If $p \ge 1$, $(1 + e_{1})(1 - e_{1}) = 1 - e_{1}^{2} = 0$ with both factors nonzero. If $p = 0$ and $q \ge 3$, $w = e_{1}e_{2}e_{3}$ has $w^{2} = (-1)^{3}e_{1}^{2}e_{2}^{2}e_{3}^{2} = (-1)(-1)^{3} = 1$ (reversing three letters costs three transpositions), so $(1 + w)(1 - w) = 0$. The remaining signatures are $(0,0)$, $(0,1)$, $(0,2)$. The Clifford conjugate $\bar{x}$ (grade involution after reversion) negates every blade of grade $1$ and $2$; for $x = \sum_{S}x_{S}e_{S}$ in these algebras $e_{S}\bar{e}_{S} = 1$ for every blade ($e_{1}(-e_{1}) = 1$, $e_{1}e_{2}(-e_{1}e_{2}) = 1$), and the cross terms cancel in pairs: for $S \ne T$, $x_{S}x_{T}(e_{S}\bar{e}_{T} + e_{T}\bar{e}_{S}) = 0$, because $\bar{e}_{T} = -e_{T}$ when $T \ne \emptyset$ and $e_{S}e_{T} = -e_{T}e_{S}$ for distinct nonempty $S, T$ in $\Cl(0,2)$, while for $S = \emptyset$ the pair is $x_{\emptyset}x_{T}(\bar{e}_{T} + e_{T}) = 0$. So $x\bar{x} = \sum_{S}x_{S}^{2}$, likewise $\bar{x}x$, which is a positive rational for $x \ne 0$; then $x^{-1} = \bar{x}/(x\bar{x})$ and $\Cl_{\Q}(p,q)$ is a division ring. The list is that of Theorem \claimref{transport.exceptional-division-algebras}. The check confirms, for every signature with $n \le 5$, the existence of nonzero null vectors exactly for $p, q \ge 1$ (on a box of coordinates), the integral zero divisors, and $x\bar{x} = \bar{x}x = \sum_{S}x_{S}^{2}$ on a grid of rational elements of $\Cl(0,0)$, $\Cl(0,1)$, $\Cl(0,2)$, the only signatures without a witness.
